\documentclass[11pt,authoryear]{elsarticle}
\usepackage[margin=1in]{geometry}
\makeatletter
\def\ps@pprintTitle{%
  \let\@oddhead\@empty
  \let\@evenhead\@empty
  \let\@oddfoot\@empty
  \let\@evenfoot\@oddfoot
}
\makeatother
\usepackage{natbib}
%\usepackage[notcite, notref]{showkeys} % show labels when printed
\usepackage{amsmath}
\usepackage{amsfonts}
\usepackage{amssymb}
\usepackage{amsthm}
\usepackage{booktabs}
\usepackage{setspace}
\usepackage{url}
\usepackage[dvipsnames,svgnames,table]{xcolor}
\usepackage{hyperref}
\hypersetup{
    colorlinks=true,
    linkcolor=blue,
    citecolor=blue,
    urlcolor=cyan
}
\onehalfspacing

\def\Real{\mathbb{R}}
\def\P#1{{\mathbb{P}}\left(#1\right)}
\def\E#1{{\mathbb E}\left[#1\right]}
\def\Var#1{{\rm Var}\left(#1\right)}
\def\cov#1{{\rm Cov}\left(#1\right)}

\newtheorem{lemma}{Lemma}
\newtheorem{theorem}{Theorem}
\newtheorem{corollary}{Corollary}
\newtheorem{remark}{Remark}
\newtheorem{proposition}{Proposition}
\newtheorem{model}{Model}
\numberwithin{equation}{section}

\begin{document}
\begin{titlepage}
\thispagestyle{empty}
\title{Binomial thinning credibility model for total and extreme claim counts}
\author{Author}
\begin{abstract}
To be written.
\end{abstract}
\maketitle
\textbf{Keywords:} credibility; binomial thinning; extreme claims;
Poisson-gamma; probit random effect

\vspace{.2cm}
Version of \today
\end{titlepage}
\newpage


\section{Introduction}
\label{sec.intro}

In many insurance portfolios, two claim counts are recorded for the same
risk unit and period: the total number of claims and the number of
extreme claims among them, for example the claims whose amounts exceed
a large-loss threshold. The two counts are nested, since the number of
extreme claims cannot exceed the total number of claims, and they are
observed repeatedly over time. A joint model for such data should
respect the nesting, accommodate time-varying covariates, and allow the
past experience of a risk unit to update the predictions of both counts.

This paper proposes a model in which the extreme claim count is a
binomial thinning of the total claim count. The total claim count of a
region follows a Poisson distribution with a gamma random effect, and
each claim is classified as extreme with a probit probability that
depends on a Gaussian random effect. Under independent and static
random effects, we show that the likelihood factorizes into a
Poisson-gamma part and a Bernoulli-probit part, so that the two sets of
parameters can be estimated separately. We further derive the posterior
distributions of the random effects and the exact Bayesian one-step
predictions of both counts in closed form. We illustrate the model with
annual claim counts of 17 Korean regions over 2015--2024.

The paper is organized as follows. Section~\ref{sec.prelim} introduces
the notation and two auxiliary results on Gaussian random effects in
probit models. Section~\ref{sec.thin} presents the model and the main
results. Section~\ref{sec.empirical} applies the model to a regional
panel of total and extreme claim counts.


\section{Preliminaries}
\label{sec.prelim}

\subsection{Notation}

Let \(\mathbb N\) denote the set of positive integers and
\(\mathbb N_0:=\mathbb N\cup\{0\}\). We write \(\Real\),
\(\Real_{\ge0}\), and \(\Real_{>0}\) for the real line, the nonnegative
real line, and the positive real line, respectively. For
\(k\in\mathbb N\), vectors in \(\Real^k\) are column vectors, and we
write \(x_{1:k}:=(x_1,\ldots,x_k)^\top\). The \(k\times k\) identity
matrix is denoted by \(\boldsymbol I_k\). Inequalities between vectors
are understood componentwise. A product over an empty index set equals
\(1\), and a sum over an empty index set equals \(0\).

We use the following distributions.
\begin{itemize}
  \setlength\itemsep{0em}
  \item \(\operatorname{Gamma}(a,b)\): gamma distribution with shape
  \(a\in\Real_{>0}\) and rate \(b\in\Real_{>0}\), so that its mean is
  \(a/b\).
  \item \(\operatorname{N}(\mu,\sigma^2)\): normal distribution with mean
  \(\mu\in\Real\) and variance \(\sigma^2\in\Real_{>0}\).
  \item \(\operatorname{MVN}_k(\mu_{1:k},\boldsymbol\Sigma)\):
  \(k\)-variate normal distribution with mean \(\mu_{1:k}\in\Real^k\) and
  positive semidefinite covariance matrix
  \(\boldsymbol\Sigma\in\Real^{k\times k}\).
  \item \(\operatorname{Pois}(\lambda)\): Poisson distribution with mean
  \(\lambda\in\Real_{>0}\).
  \item \(\operatorname{Bin}(n,p)\): binomial distribution with
  \(n\in\mathbb N_0\) trials and success probability \(p\in(0,1)\), where
  \(\operatorname{Bin}(0,p)\) is the point mass at \(0\).
\end{itemize}

Let \(\phi\) and \(\Phi\) denote the density and the distribution
function of \(\operatorname{N}(0,1)\). For a positive semidefinite
matrix \(\boldsymbol\Sigma\in\Real^{k\times k}\), let
\[
\Phi_k(x_{1:k};\boldsymbol\Sigma)
:=
\P{X_{1:k}\le x_{1:k}},
\qquad
X_{1:k}\sim\operatorname{MVN}_k(0_{1:k},\boldsymbol\Sigma),
\quad
x_{1:k}\in\Real^k,
\]
where \(0_{1:k}\in\Real^k\) is the zero vector. For \(k=0\), we set
\(\Phi_0:=1\).


\subsection{Gaussian random effects in probit models}

The following identity converts the expectation of a product of probit
functions of a Gaussian vector into a multivariate normal distribution
function.

\begin{lemma}[Gaussian product-probit identity]
\label{lem.probit}
Let \(X_{1:k}\sim\operatorname{MVN}_k(\mu_{1:k},\boldsymbol\Sigma)\),
where \(\mu_{1:k}\in\Real^k\) and \(\boldsymbol\Sigma\in\Real^{k\times k}\)
is positive semidefinite. Then
\[
\E{\prod_{j=1}^k\Phi(X_j)}
=
\Phi_k\left(\mu_{1:k};\boldsymbol I_k+\boldsymbol\Sigma\right).
\]
\end{lemma}

\begin{proof}
Let \(W_{1:k}\sim\operatorname{MVN}_k(0_{1:k},\boldsymbol I_k)\) be
independent of \(X_{1:k}\). Since \(W_1,\ldots,W_k\) are independent
standard normal random variables,
\[
\prod_{j=1}^k\Phi(X_j)
=
\P{W_{1:k}\le X_{1:k}\mid X_{1:k}}.
\]
Taking expectations gives
\[
\E{\prod_{j=1}^k\Phi(X_j)}
=
\P{W_{1:k}\le X_{1:k}}
=
\P{W_{1:k}-\left(X_{1:k}-\mu_{1:k}\right)\le\mu_{1:k}}.
\]
Since \(W_{1:k}\) and \(X_{1:k}\) are independent Gaussian vectors,
\[
W_{1:k}-\left(X_{1:k}-\mu_{1:k}\right)
\sim
\operatorname{MVN}_k\left(0_{1:k},\boldsymbol I_k+\boldsymbol\Sigma\right),
\]
which proves the result.
\end{proof}

Lemma~\ref{lem.probit} does not require \(\boldsymbol\Sigma\) to be
positive definite. This is essential for the next result, in which a
single scalar random effect is shared by all components, so that the
covariance matrix has rank one.

\begin{lemma}[Bernoulli-probit random effect]
\label{lem.bp}
Let \(\Theta\sim\operatorname{N}(0,\sigma^2)\) with \(\sigma>0\), and let
\(k\in\mathbb N\), \(u_{1:k}\in\{-1,1\}^k\), and \(h_{1:k}\in\Real^k\).
Then
\[
\begin{aligned}
\E{\prod_{j=1}^k\Phi\left(u_j(h_j+\Theta)\right)}
&=
\int_{\Real}\phi(z)\prod_{j=1}^k\Phi\left(u_j(h_j+\sigma z)\right)
{\rm d}z\\
&=
\Phi_k\left(
(u_1h_1,\ldots,u_kh_k)^\top;
\boldsymbol I_k+\sigma^2u_{1:k}u_{1:k}^\top
\right).
\end{aligned}
\]
\end{lemma}

\begin{proof}
The first equality follows since \(\Theta\) has the same distribution
as \(\sigma Z\) with \(Z\sim\operatorname{N}(0,1)\). For the second equality, write
\[
X_{1:k}
:=
\left(u_1(h_1+\Theta),\ldots,u_k(h_k+\Theta)\right)^\top
=
(u_1h_1,\ldots,u_kh_k)^\top+u_{1:k}\Theta.
\]
Then
\(
X_{1:k}\sim\operatorname{MVN}_k\left(
(u_1h_1,\ldots,u_kh_k)^\top,\sigma^2u_{1:k}u_{1:k}^\top
\right),
\)
and Lemma~\ref{lem.probit} proves the result.
\end{proof}

\begin{remark}\label{rem.bp}
Lemma~\ref{lem.bp} has the following interpretation. Let
\(Z_1,\ldots,Z_k\) be conditionally independent given \(\Theta\) with
\(Z_j\mid\Theta\sim\operatorname{Bin}\left(1,\Phi(h_j+\Theta)\right)\).
Since \(1-\Phi(a)=\Phi(-a)\) for \(a\in\Real\), we have
\(
\P{Z_j=z_j\mid\Theta}=\Phi\left(u_j(h_j+\Theta)\right)
\)
with \(u_j:=2z_j-1\) for \(z_j\in\{0,1\}\). Hence
Lemma~\ref{lem.bp} gives the joint probability
\(\P{Z_{1:k}=z_{1:k}}\) in closed form. The covariance matrix
\(\sigma^2u_{1:k}u_{1:k}^\top\) is of rank one, reflecting the single
random effect shared by all trials, and the one-dimensional integral
representation is convenient for computation when \(k\) is large.
\end{remark}


%==========================================================================
% Section: Poisson-gamma / binomial-probit thinning model
%          for bivariate claim frequency
%==========================================================================
\section{Binomial thinning model for bivariate claim frequency}
\label{sec.thin}

In this section, we develop a model for two nested claim counts.
For each region \(i\) and period \(t\), let \(N_{it}\) denote the total
number of claims and let \(M_{it}\) denote the number of extreme claims
among them, for example, the claims whose amounts exceed a given
threshold. By construction,
\[
0\le M_{it}\le N_{it}.
\]
A natural way to respect this constraint is to regard \(M_{it}\) as a
binomial thinning of \(N_{it}\): each of the \(N_{it}\) claims is
classified as extreme or not. Repeated observations of the same region
over time are linked through two region-specific random effects, one
for the claim frequency and one for the probability that a claim is
extreme. Throughout this section, the two random effects are assumed to
be independent and static over time, and predictions are obtained as
exact Bayesian posterior means.

\begin{model}[Poisson-gamma binomial-probit thinning model]
\label{model.thin}
Let
\[
r,\sigma\in\Real_{>0},
\qquad
\alpha,\beta\in\Real^p,
\]
and let \(x_{it}\in\Real^p\), \(i=1,\ldots,I\), \(t\in\mathbb N\), be
predetermined exogenous covariates, which are allowed to vary over time.
Define
\[
\lambda_{it}:=\exp\left(x_{it}^\top\alpha\right)\in\Real_{>0},
\qquad
\eta_{it}:=x_{it}^\top\beta\in\Real.
\]
We define the bivariate sequences
\(\left((N_{it},M_{it})\right)_{t\in\mathbb N}\), \(i=1,\ldots,I\),
through the following random effect assumption and observation
equations.
\begin{itemize}
  \item[i.] Random effects: For each region \(i\), assume
  \[
  \Theta_i^{(1)}\sim\operatorname{Gamma}(r,r),
  \qquad
  \Theta_i^{(2)}\sim\operatorname{N}(0,\sigma^2),
  \]
  where \(\Theta_i^{(1)}\) and \(\Theta_i^{(2)}\) are independent.
  Furthermore, the pairs
  \(\left(\Theta_i^{(1)},\Theta_i^{(2)}\right)\), \(i=1,\ldots,I\),
  are mutually independent.

  \item[ii.] Observation equations: Conditional on
  \(\left(\Theta_i^{(1)},\Theta_i^{(2)}\right)\), assume that the pairs
  \((N_{it},M_{it})\), \(t\in\mathbb N\), are mutually independent with
  \[
  N_{it}\mid\Theta_i^{(1)},\Theta_i^{(2)}
  \sim
  \operatorname{Pois}\left(\lambda_{it}\Theta_i^{(1)}\right)
  \]
  and
  \[
  M_{it}\mid N_{it},\Theta_i^{(1)},\Theta_i^{(2)}
  \sim
  \operatorname{Bin}\left(N_{it},\Phi\left(\eta_{it}+\Theta_i^{(2)}\right)\right),
  \]
  Furthermore, the regions are independent: the sequences
  \(\left((N_{it},M_{it})\right)_{t\in\mathbb N}\), \(i=1,\ldots,I\),
  are mutually independent.
\end{itemize}
\end{model}
Here, \(\operatorname{Bin}(0,p)\) is the point mass at \(0\), so that
\(M_{it}=0\) whenever \(N_{it}=0\).

Since \(\E{\Theta_i^{(1)}}=1\) and \(\E{\Theta_i^{(2)}}=0\), the
intercepts in \(\alpha\) and \(\beta\) are identifiable, and
\(\lambda_{it}\) is the a priori expected total claim count. The random
effect \(\Theta_i^{(1)}\) captures the unobserved heterogeneity of the
claim frequency across regions, whereas \(\Theta_i^{(2)}\) captures the
unobserved heterogeneity of the propensity of a claim to be extreme.
For notational convenience, write
\[
p_{it}(\theta):=\Phi\left(\eta_{it}+\theta\right),
\qquad \theta\in\Real.
\]

\begin{remark}\label{rem.thin.exposure}
All results in this section use \(\lambda_{it}\) and \(\eta_{it}\) only
as known constants. Hence they remain valid for any
\(\lambda_{it}\in\Real_{>0}\) and \(\eta_{it}\in\Real\); in particular,
an exposure offset \(e_{it}>0\) can be included through
\(\lambda_{it}=e_{it}\exp\left(x_{it}^\top\alpha\right)\).
\end{remark}

The following lemma is the well-known thinning property of the Poisson
distribution. It shows that, conditional on the random effects, the
extreme and non-extreme claim counts are independent Poisson random
variables.

\begin{lemma}[Poisson thinning]
\label{lem.thin}
Under Model~\ref{model.thin}, conditional on
\(\left(\Theta_i^{(1)},\Theta_i^{(2)}\right)\), the random variables
\(M_{it}\) and \(N_{it}-M_{it}\) are independent with
\[
M_{it}\mid\Theta_i^{(1)},\Theta_i^{(2)}
\sim
\operatorname{Pois}\left(
\lambda_{it}\Theta_i^{(1)}p_{it}\left(\Theta_i^{(2)}\right)
\right)
\]
and
\[
N_{it}-M_{it}\mid\Theta_i^{(1)},\Theta_i^{(2)}
\sim
\operatorname{Pois}\left(
\lambda_{it}\Theta_i^{(1)}
\left\{1-p_{it}\left(\Theta_i^{(2)}\right)\right\}
\right).
\]
\end{lemma}

\begin{proof}
Fix \(\theta_1>0\) and \(\theta_2\in\Real\), and write
\(\Lambda:=\lambda_{it}\theta_1\) and \(p:=p_{it}(\theta_2)\). For
\(m,k\in\mathbb N_0\), Model~\ref{model.thin} gives
\[
\begin{aligned}
&\P{M_{it}=m,\,N_{it}-M_{it}=k\mid
\Theta_i^{(1)}=\theta_1,\Theta_i^{(2)}=\theta_2}\\
&\quad=
\P{N_{it}=m+k\mid\Theta_i^{(1)}=\theta_1}
\binom{m+k}{m}p^m(1-p)^k\\
&\quad=
\frac{e^{-\Lambda}\Lambda^{m+k}}{(m+k)!}
\frac{(m+k)!}{m!\,k!}p^m(1-p)^k\\
&\quad=
\frac{e^{-\Lambda p}(\Lambda p)^m}{m!}
\cdot
\frac{e^{-\Lambda(1-p)}\left\{\Lambda(1-p)\right\}^k}{k!},
\end{aligned}
\]
which proves the result.
\end{proof}

Since the regions are independent, it suffices to study a single
region. The following remark collects the notation for its observed
history.

\begin{remark}\label{rem.thin.notation}
Fix a region \(i\) and \(t\in\mathbb N\), and consider a realized
history
\[
\mathcal H_{it}:=
\left\{
(N_{ij},M_{ij})=(n_{ij},m_{ij}):\ j=1,\ldots,t
\right\},
\qquad
n_{ij},m_{ij}\in\mathbb N_0,
\quad
m_{ij}\le n_{ij}.
\]
Define the cumulative claim count and the cumulative a priori
frequency by
\[
n_{i\bullet t}:=\sum_{j=1}^t n_{ij},
\qquad
\lambda_{i\bullet t}:=\sum_{j=1}^t \lambda_{ij}.
\]

Each of the \(n_{i\bullet t}\) observed claims is regarded as a
Bernoulli trial, which is a success if the claim is extreme. To this
end, for \(j=1,\ldots,t\), define the sign vector
\[
u^{(j)}:=
\big(
\underbrace{1,\ldots,1}_{m_{ij}},
\underbrace{-1,\ldots,-1}_{n_{ij}-m_{ij}}
\big)^\top
\in\{-1,1\}^{n_{ij}},
\]
which is empty when \(n_{ij}=0\), and stack these vectors as
\[
u_{1:n_{i\bullet t}}:=
\left(
u^{(1)\top},\ldots,u^{(t)\top}
\right)^\top,
\qquad
h_{1:n_{i\bullet t}}:=
\big(
\underbrace{\eta_{i1},\ldots,\eta_{i1}}_{n_{i1}},
\ldots,
\underbrace{\eta_{it},\ldots,\eta_{it}}_{n_{it}}
\big)^\top .
\]
That is, \(u_k\) is the sign of the \(k\)th claim and \(h_k\) is the
value of \(\eta_{ij}\) in the period \(j\) of the \(k\)th claim. We then
define
\[
\mu_{it}:=
\left(u_1h_1,\ldots,u_{n_{i\bullet t}}h_{n_{i\bullet t}}\right)^\top
\in\Real^{n_{i\bullet t}},
\qquad
\boldsymbol U_{it}:=
u_{1:n_{i\bullet t}}u_{1:n_{i\bullet t}}^\top
\in\Real^{n_{i\bullet t}\times n_{i\bullet t}},
\]
which are the mean vector and the rank-one matrix appearing in
Lemma~\ref{lem.bp} for the trials formed by the observed claims.

For \(n,m\in\mathbb N_0\) with \(m\le n\), let
\(\mu_{it}^{(n,m)}\in\Real^{n_{i\bullet t}+n}\) and
\(\boldsymbol U_{it}^{(n,m)}\in\Real^{(n_{i\bullet t}+n)\times(n_{i\bullet t}+n)}\)
denote the corresponding quantities for the augmented history
\(\mathcal H_{it}\cup\{(N_{i,t+1},M_{i,t+1})=(n,m)\}\). That is,
\(n\) additional trials with \(\eta\)-value \(\eta_{i,t+1}\) are appended,
\(m\) of them with sign \(1\) and \(n-m\) of them with sign \(-1\).

Finally, define
\[
D_{it}(\theta):=
\prod_{j=1}^t
p_{ij}(\theta)^{m_{ij}}
\left\{1-p_{ij}(\theta)\right\}^{n_{ij}-m_{ij}},
\qquad\theta\in\Real.
\]
Since \(1-\Phi(a)=\Phi(-a)\) for \(a\in\Real\), we have
\begin{equation}\label{eq.thin.D}
D_{it}(\theta)
=
\prod_{k=1}^{n_{i\bullet t}}
\Phi\left(u_k(h_k+\theta)\right).
\end{equation}
For \(n_{i\bullet t}=0\), we adopt the conventions
\(D_{it}(\theta)=1\) and
\(\Phi_0\left(\mu_{it};\boldsymbol I_0+\sigma^2\boldsymbol U_{it}\right)=1\).
\end{remark}

The following theorem shows that the joint probability of the observed
history factorizes into a closed-form Poisson-gamma part, which depends
only on \((\alpha,r)\), and a Bernoulli-probit part of the form in
Lemma~\ref{lem.bp}, which depends only on \((\beta,\sigma)\).

\begin{theorem}[Likelihood factorization]
\label{thm.thin.lik}
Consider Model~\ref{model.thin} with the realized history
\(\mathcal H_{it}\), and let the notation be as in
Remark~\ref{rem.thin.notation}. Then
\begin{equation}\label{eq.thin.factorization}
\P{\mathcal H_{it}}
=
L^{(1)}_{it}\,L^{(2)}_{it},
\end{equation}
where
\begin{equation}\label{eq.thin.L1}
L^{(1)}_{it}
:=
\P{N_{i,1:t}=n_{i,1:t}}
=
\frac{\Gamma\left(r+n_{i\bullet t}\right)}
{\Gamma(r)\prod_{j=1}^t n_{ij}!}
\,
\frac{r^r\prod_{j=1}^t\lambda_{ij}^{n_{ij}}}
{\left(r+\lambda_{i\bullet t}\right)^{r+n_{i\bullet t}}}
\end{equation}
and
\begin{equation}\label{eq.thin.L2}
\begin{aligned}
L^{(2)}_{it}
&:=
\P{M_{i,1:t}=m_{i,1:t}\mid N_{i,1:t}=n_{i,1:t}}\\
&=
\prod_{j=1}^t\binom{n_{ij}}{m_{ij}}
\int_{\Real}
\frac{1}{\sigma}\phi\left(\frac{\theta}{\sigma}\right)
D_{it}(\theta)
\,{\rm d}\theta\\
&=
\prod_{j=1}^t\binom{n_{ij}}{m_{ij}}
\,
\Phi_{n_{i\bullet t}}
\left(
\mu_{it};
\boldsymbol I_{n_{i\bullet t}}
+
\sigma^2\boldsymbol U_{it}
\right).
\end{aligned}
\end{equation}
\end{theorem}

\begin{proof}
By the conditional independence in Model~\ref{model.thin},
\[
\P{\mathcal H_{it}}
=
\E{
\prod_{j=1}^t
\P{N_{ij}=n_{ij}\mid\Theta_i^{(1)}}
\P{M_{ij}=m_{ij}\mid N_{ij}=n_{ij},\Theta_i^{(2)}}
}.
\]
The first factor in the product depends only on \(\Theta_i^{(1)}\) and
the second only on \(\Theta_i^{(2)}\). Since \(\Theta_i^{(1)}\) and
\(\Theta_i^{(2)}\) are independent, we obtain
\[
\P{\mathcal H_{it}}
=
\E{
\prod_{j=1}^t
\P{N_{ij}=n_{ij}\mid\Theta_i^{(1)}}
}
\,
\E{
\prod_{j=1}^t
\P{M_{ij}=m_{ij}\mid N_{ij}=n_{ij},\Theta_i^{(2)}}
}.
\]

For the first expectation, the gamma integral gives
\[
\begin{aligned}
\E{
\prod_{j=1}^t
\P{N_{ij}=n_{ij}\mid\Theta_i^{(1)}}
}
&=
\int_0^\infty
\prod_{j=1}^t
\frac{e^{-\lambda_{ij}\theta}(\lambda_{ij}\theta)^{n_{ij}}}{n_{ij}!}
\,
\frac{r^r}{\Gamma(r)}\theta^{r-1}e^{-r\theta}
\,{\rm d}\theta\\
&=
\frac{r^r\prod_{j=1}^t\lambda_{ij}^{n_{ij}}}
{\Gamma(r)\prod_{j=1}^t n_{ij}!}
\int_0^\infty
\theta^{r+n_{i\bullet t}-1}
e^{-(r+\lambda_{i\bullet t})\theta}
\,{\rm d}\theta,
\end{aligned}
\]
which equals the right-hand side of \eqref{eq.thin.L1}. Summing the
first display over all \(m_{i,1:t}\) with \(0\le m_{ij}\le n_{ij}\)
shows that this expectation is \(\P{N_{i,1:t}=n_{i,1:t}}\), because
the second expectation sums to one.

For the second expectation, we have
\[
\E{
\prod_{j=1}^t
\P{M_{ij}=m_{ij}\mid N_{ij}=n_{ij},\Theta_i^{(2)}}
}
=
\prod_{j=1}^t\binom{n_{ij}}{m_{ij}}
\E{D_{it}\left(\Theta_i^{(2)}\right)},
\]
which proves the second line of \eqref{eq.thin.L2}. It then follows
from \eqref{eq.thin.factorization} and \eqref{eq.thin.L1} that this
expression is \(\P{M_{i,1:t}=m_{i,1:t}\mid N_{i,1:t}=n_{i,1:t}}\).
Finally, by \eqref{eq.thin.D},
\[
\E{D_{it}\left(\Theta_i^{(2)}\right)}
=
\E{
\prod_{k=1}^{n_{i\bullet t}}
\Phi\left(u_k\left(h_k+\Theta_i^{(2)}\right)\right)
},
\]
and applying Lemma~\ref{lem.bp} with \(k=n_{i\bullet t}\),
\(u_{1:k}\), and \(h_{1:k}\) proves the last line of
\eqref{eq.thin.L2}.
\end{proof}

\begin{remark}[Separable estimation]
\label{rem.thin.estimation}
By Theorem~\ref{thm.thin.lik} and the independence of the regions,
the log-likelihood of the observations
\(\mathcal H_{1T},\ldots,\mathcal H_{IT}\) decomposes as
\[
\ell(\alpha,r,\beta,\sigma)
=
\ell^{(1)}(\alpha,r)+\ell^{(2)}(\beta,\sigma),
\qquad
\ell^{(k)}:=\sum_{i=1}^I\log L^{(k)}_{iT},
\quad k=1,2.
\]
Hence the maximum likelihood estimators of \((\alpha,r)\) and
\((\beta,\sigma)\) are obtained from two separate optimizations, and
the Fisher information matrix is block-diagonal. Here,
\(\ell^{(1)}\) is the log-likelihood of a negative binomial
random effect panel model for \(N_{it}\), and \(\ell^{(2)}\) is the
log-likelihood of a binomial probit model for \(M_{it}\) given
\(N_{it}\) with a Gaussian random intercept. Periods with \(n_{it}=0\)
do not contribute to \(\ell^{(2)}\).
\end{remark}

\begin{remark}[Computation]
\label{rem.thin.computation}
The dimension \(n_{i\bullet t}\) of the multivariate normal
distribution function in \eqref{eq.thin.L2} equals the cumulative
number of claims of region \(i\), which can be large. However, since
\(\boldsymbol U_{it}\) is of rank one, the second line of
\eqref{eq.thin.L2} expresses the same quantity as the one-dimensional
integral
\[
\Phi_{n_{i\bullet t}}
\left(
\mu_{it};
\boldsymbol I_{n_{i\bullet t}}+\sigma^2\boldsymbol U_{it}
\right)
=
\int_{\Real}\phi(z)\,D_{it}(\sigma z)\,{\rm d}z,
\]
which depends on the data only through the \(t\) pairs
\((n_{ij},m_{ij})\) and can be evaluated accurately by Gauss--Hermite
quadrature.
\end{remark}

The following corollary shows that the two random effects remain
independent a posteriori, and that the frequency random effect has a
conjugate gamma posterior.

\begin{corollary}[Posterior distribution of the random effects]
\label{cor.thin.posterior}
Under the setting of Theorem~\ref{thm.thin.lik}, the following
results hold.
\begin{itemize}
  \item[i.] Conditional on \(\mathcal H_{it}\), the random effects
  \(\Theta_i^{(1)}\) and \(\Theta_i^{(2)}\) are independent.

  \item[ii.] We have
  \[
  \Theta_i^{(1)}\mid\mathcal H_{it}
  \sim
  \operatorname{Gamma}\left(r+n_{i\bullet t},\,r+\lambda_{i\bullet t}\right).
  \]

  \item[iii.] The conditional distribution
  \(\Theta_i^{(2)}\mid\mathcal H_{it}\) has density
  \(\pi(\theta\mid\mathcal H_{it})\) at \(\theta\in\Real\) given by
  \begin{equation}\label{eq.thin.posterior2}
  \pi(\theta\mid\mathcal H_{it})
  =
  \frac{
  \displaystyle
  \frac{1}{\sigma}\phi\left(\frac{\theta}{\sigma}\right)
  \prod_{k=1}^{n_{i\bullet t}}
  \Phi\left(u_k(h_k+\theta)\right)
  }{
  \Phi_{n_{i\bullet t}}
  \left(
  \mu_{it};
  \boldsymbol I_{n_{i\bullet t}}+\sigma^2\boldsymbol U_{it}
  \right)
  }.
  \end{equation}
\end{itemize}
\end{corollary}

\begin{proof}
Let \(\pi_1\) and \(\pi_2\) denote the prior densities of
\(\Theta_i^{(1)}\) and \(\Theta_i^{(2)}\), respectively. By Bayes'
formula and the proof of Theorem~\ref{thm.thin.lik}, the joint
posterior density of
\(\left(\Theta_i^{(1)},\Theta_i^{(2)}\right)\) at
\((\theta_1,\theta_2)\in\Real_{>0}\times\Real\) is
\[
\frac{
\pi_1(\theta_1)
\prod_{j=1}^t\P{N_{ij}=n_{ij}\mid\Theta_i^{(1)}=\theta_1}
}{L^{(1)}_{it}}
\cdot
\frac{
\pi_2(\theta_2)
\prod_{j=1}^t\binom{n_{ij}}{m_{ij}}D_{it}(\theta_2)
}{L^{(2)}_{it}}.
\]
Each factor integrates to one, which proves part i. The first factor
is proportional to
\[
\theta_1^{r+n_{i\bullet t}-1}e^{-(r+\lambda_{i\bullet t})\theta_1},
\]
which proves part ii. Substituting \eqref{eq.thin.D} and
\eqref{eq.thin.L2} into the second factor proves part iii.
\end{proof}

By part ii of Corollary~\ref{cor.thin.posterior}, the posterior of the
frequency random effect depends on the history only through
\(n_{i\bullet t}\); in particular, the extreme claim counts
\(m_{i,1:t}\) carry no information on \(\Theta_i^{(1)}\). Conversely,
by part iii, the total claim counts affect the posterior of
\(\Theta_i^{(2)}\) only by determining the number of Bernoulli trials.
The density \eqref{eq.thin.posterior2} is a Gaussian density
multiplied by a product of probit functions, normalized by the
multivariate normal distribution function in Lemma~\ref{lem.bp}.

We now present the exact Bayesian one-step-ahead predictions.

\begin{theorem}[Closed-form one-step prediction]
\label{thm.thin.pred}
Consider Model~\ref{model.thin} with the realized history
\(\mathcal H_{it}\), and let the notation be as in
Remark~\ref{rem.thin.notation}. Then, we have the following results.
\begin{itemize}
  \item[i.] For \(n,m\in\mathbb N_0\) with \(m\le n\), the predictive
  probability is
  \[
  \begin{aligned}
  &\P{N_{i,t+1}=n,\,M_{i,t+1}=m\mid\mathcal H_{it}}\\
  &\quad=
  \frac{\Gamma\left(r+n_{i\bullet t}+n\right)}
  {\Gamma\left(r+n_{i\bullet t}\right)n!}
  \left(
  \frac{r+\lambda_{i\bullet t}}{r+\lambda_{i\bullet t}+\lambda_{i,t+1}}
  \right)^{r+n_{i\bullet t}}
  \left(
  \frac{\lambda_{i,t+1}}{r+\lambda_{i\bullet t}+\lambda_{i,t+1}}
  \right)^{n}\\
  &\qquad\times
  \binom{n}{m}
  \frac{
  \Phi_{n_{i\bullet t}+n}
  \left(
  \mu_{it}^{(n,m)};
  \boldsymbol I_{n_{i\bullet t}+n}+\sigma^2\boldsymbol U_{it}^{(n,m)}
  \right)
  }{
  \Phi_{n_{i\bullet t}}
  \left(
  \mu_{it};
  \boldsymbol I_{n_{i\bullet t}}+\sigma^2\boldsymbol U_{it}
  \right)
  }.
  \end{aligned}
  \]

  \item[ii.] The one-step-ahead prediction of the total claim count is
  \[
  \E{N_{i,t+1}\mid\mathcal H_{it}}
  =
  \lambda_{i,t+1}\,
  \frac{r+n_{i\bullet t}}{r+\lambda_{i\bullet t}}
  =
  \lambda_{i,t+1}
  \left\{
  \zeta_{it}\frac{n_{i\bullet t}}{\lambda_{i\bullet t}}
  +
  \left(1-\zeta_{it}\right)
  \right\},
  \qquad
  \zeta_{it}:=\frac{\lambda_{i\bullet t}}{r+\lambda_{i\bullet t}}.
  \]

  \item[iii.] The predictive probability that a claim in period
  \(t+1\) is extreme is
  \[
  \widehat p_{i,t+1}
  :=
  \E{p_{i,t+1}\left(\Theta_i^{(2)}\right)\mid\mathcal H_{it}}
  =
  \frac{
  \Phi_{n_{i\bullet t}+1}
  \left(
  \mu_{it}^{(1,1)};
  \boldsymbol I_{n_{i\bullet t}+1}+\sigma^2\boldsymbol U_{it}^{(1,1)}
  \right)
  }{
  \Phi_{n_{i\bullet t}}
  \left(
  \mu_{it};
  \boldsymbol I_{n_{i\bullet t}}+\sigma^2\boldsymbol U_{it}
  \right)
  }.
  \]

  \item[iv.] For \(n\in\mathbb N_0\), the one-step-ahead prediction of
  the extreme claim count given the total claim count is
  \[
  \E{M_{i,t+1}\mid\mathcal H_{it},N_{i,t+1}=n}
  =
  n\,\widehat p_{i,t+1}.
  \]

  \item[v.] The one-step-ahead prediction of the extreme claim count is
  \[
  \E{M_{i,t+1}\mid\mathcal H_{it}}
  =
  \lambda_{i,t+1}\,
  \frac{r+n_{i\bullet t}}{r+\lambda_{i\bullet t}}\,
  \widehat p_{i,t+1}.
  \]
\end{itemize}
\end{theorem}

\begin{proof}
We first show part i. Applying Theorem~\ref{thm.thin.lik} to the
augmented history
\(\mathcal H_{it}\cup\{(N_{i,t+1},M_{i,t+1})=(n,m)\}\) and to
\(\mathcal H_{it}\), we obtain
\[
\P{N_{i,t+1}=n,\,M_{i,t+1}=m\mid\mathcal H_{it}}
=
\frac{L^{(1)}_{i,t+1}}{L^{(1)}_{it}}
\cdot
\frac{L^{(2)}_{i,t+1}}{L^{(2)}_{it}},
\]
where \(L^{(1)}_{i,t+1}\) and \(L^{(2)}_{i,t+1}\) are evaluated at the
augmented history. By \eqref{eq.thin.L1},
\[
\frac{L^{(1)}_{i,t+1}}{L^{(1)}_{it}}
=
\frac{\Gamma\left(r+n_{i\bullet t}+n\right)}
{\Gamma\left(r+n_{i\bullet t}\right)n!}
\,
\frac{\lambda_{i,t+1}^{n}\left(r+\lambda_{i\bullet t}\right)^{r+n_{i\bullet t}}}
{\left(r+\lambda_{i\bullet t}+\lambda_{i,t+1}\right)^{r+n_{i\bullet t}+n}},
\]
and, by the last line of \eqref{eq.thin.L2},
\[
\frac{L^{(2)}_{i,t+1}}{L^{(2)}_{it}}
=
\binom{n}{m}
\frac{
\Phi_{n_{i\bullet t}+n}
\left(
\mu_{it}^{(n,m)};
\boldsymbol I_{n_{i\bullet t}+n}+\sigma^2\boldsymbol U_{it}^{(n,m)}
\right)
}{
\Phi_{n_{i\bullet t}}
\left(
\mu_{it};
\boldsymbol I_{n_{i\bullet t}}+\sigma^2\boldsymbol U_{it}
\right)
},
\]
which proves part i.

To show part ii, note that, conditional on
\(\left(\Theta_i^{(1)},\Theta_i^{(2)}\right)\), \(N_{i,t+1}\) is
independent of \(\mathcal H_{it}\). Hence, by part ii of
Corollary~\ref{cor.thin.posterior},
\[
\E{N_{i,t+1}\mid\mathcal H_{it}}
=
\lambda_{i,t+1}\E{\Theta_i^{(1)}\mid\mathcal H_{it}}
=
\lambda_{i,t+1}\frac{r+n_{i\bullet t}}{r+\lambda_{i\bullet t}},
\]
and the second expression follows from
\(
(r+n_{i\bullet t})/(r+\lambda_{i\bullet t})
=
\zeta_{it}\,n_{i\bullet t}/\lambda_{i\bullet t}+1-\zeta_{it}.
\)

To show part iii, part iii of Corollary~\ref{cor.thin.posterior} gives
\[
\widehat p_{i,t+1}
=
\frac{
\E{
p_{i,t+1}\left(\Theta_i^{(2)}\right)
\prod_{k=1}^{n_{i\bullet t}}
\Phi\left(u_k\left(h_k+\Theta_i^{(2)}\right)\right)
}
}{
\Phi_{n_{i\bullet t}}
\left(
\mu_{it};
\boldsymbol I_{n_{i\bullet t}}+\sigma^2\boldsymbol U_{it}
\right)
}.
\]
Since
\(p_{i,t+1}(\theta)=\Phi\left(1\cdot(\eta_{i,t+1}+\theta)\right)\), the
product in the numerator is the product in \eqref{eq.thin.D} for the
history augmented by \((n,m)=(1,1)\). Applying Lemma~\ref{lem.bp} to the
augmented history proves part iii.

To show part iv, we first observe that, conditional on
\(\mathcal H_{it}\), the joint density of
\(\left(\Theta_i^{(1)},\Theta_i^{(2)}\right)\) and the probability mass
function of \(N_{i,t+1}\) satisfy
\[
\pi(\theta_1,\theta_2,n\mid\mathcal H_{it})
=
\pi(\theta_1\mid\mathcal H_{it})
\P{N_{i,t+1}=n\mid\Theta_i^{(1)}=\theta_1}
\,
\pi(\theta_2\mid\mathcal H_{it}),
\]
by part i of Corollary~\ref{cor.thin.posterior} and
Model~\ref{model.thin}. Hence \(\Theta_i^{(2)}\) is independent of
\(N_{i,t+1}\) conditional on \(\mathcal H_{it}\). Together with the
binomial observation equation, this gives
\[
\E{M_{i,t+1}\mid\mathcal H_{it},N_{i,t+1}=n}
=
n\,\E{p_{i,t+1}\left(\Theta_i^{(2)}\right)\mid\mathcal H_{it},N_{i,t+1}=n}
=
n\,\widehat p_{i,t+1}.
\]

Finally, part v follows from part iv and part ii, since
\[
\E{M_{i,t+1}\mid\mathcal H_{it}}
=
\E{N_{i,t+1}\mid\mathcal H_{it}}\,\widehat p_{i,t+1}.
\]
\end{proof}

\begin{remark}[Interpretation of the predictions]
\label{rem.thin.interpretation}
Part ii of Theorem~\ref{thm.thin.pred} is the exact credibility formula
of the Poisson-gamma model: the a priori frequency \(\lambda_{i,t+1}\)
is multiplied by a credibility-weighted average of the observed ratio
\(n_{i\bullet t}/\lambda_{i\bullet t}\) and the prior mean \(1\), with
credibility factor \(\zeta_{it}\), which increases to one as
\(\lambda_{i\bullet t}\to\infty\). Part iii is the posterior predictive
probability that an additional Bernoulli-probit trial, namely a new
claim, is a success.
In particular, if \(n_{i\bullet t}=0\), then
\(
\widehat p_{i,t+1}
=
\Phi\left(\eta_{i,t+1}/\sqrt{1+\sigma^2}\right),
\)
which is the a priori probability that a claim is extreme. Part v shows
that the prediction of the extreme claim count is the product of the
two exact Bayesian predictions, one for the frequency and one for the
proportion of extreme claims, as a consequence of the a posteriori
independence of the random effects.
\end{remark}

We conclude this section with the unconditional moments of
Model~\ref{model.thin}, which describe the dependence induced by the
random effects.

\begin{proposition}[Moments]
\label{prop.thin.moments}
Under Model~\ref{model.thin}, fix a region \(i\), and define
\[
\pi_{it}:=\Phi\left(\frac{\eta_{it}}{\sqrt{1+\sigma^2}}\right),
\qquad
\pi_{i,st}:=
\Phi_2\left(
(\eta_{is},\eta_{it})^\top;
\boldsymbol I_2+\sigma^2\boldsymbol J_2
\right),
\qquad s,t\in\mathbb N,
\]
where \(\boldsymbol J_2\) is the \(2\times2\) matrix of ones. Then, for
\(s,t\in\mathbb N\) with \(s\neq t\), we have
\[
\begin{aligned}
\E{N_{it}}&=\lambda_{it},
&
\E{M_{it}}&=\lambda_{it}\pi_{it},\\
\Var{N_{it}}&=\lambda_{it}+\frac{\lambda_{it}^2}{r},
&
\cov{N_{is},N_{it}}&=\frac{\lambda_{is}\lambda_{it}}{r},\\
\cov{N_{it},M_{it}}&=\pi_{it}\left(\lambda_{it}+\frac{\lambda_{it}^2}{r}\right),
&
\cov{N_{is},M_{it}}&=\frac{\lambda_{is}\lambda_{it}\pi_{it}}{r},
\end{aligned}
\]
and
\[
\begin{aligned}
\Var{M_{it}}
&=
\lambda_{it}\pi_{it}
+
\lambda_{it}^2
\left\{
\left(1+\frac1r\right)\pi_{i,tt}-\pi_{it}^2
\right\},\\
\cov{M_{is},M_{it}}
&=
\lambda_{is}\lambda_{it}
\left\{
\left(1+\frac1r\right)\pi_{i,st}-\pi_{is}\pi_{it}
\right\}.
\end{aligned}
\]
\end{proposition}

\begin{proof}
Write \(\Theta_1:=\Theta_i^{(1)}\) and \(\Theta_2:=\Theta_i^{(2)}\), and
recall \(\E{\Theta_1}=1\) and \(\E{\Theta_1^2}=1+1/r\). By
Lemma~\ref{lem.bp} with \(k=1\) and \(k=2\), and \(u_{1:k}\) the
vector of ones,
\[
\E{p_{it}(\Theta_2)}=\pi_{it},
\qquad
\E{p_{is}(\Theta_2)p_{it}(\Theta_2)}=\pi_{i,st},
\qquad s,t\in\mathbb N.
\]
Conditional on \((\Theta_1,\Theta_2)\), Model~\ref{model.thin} and
Lemma~\ref{lem.thin} give, for \(s\neq t\),
\[
\begin{aligned}
\E{N_{it}\mid\Theta_1,\Theta_2}&=\lambda_{it}\Theta_1,
&
\E{M_{it}\mid\Theta_1,\Theta_2}&=\lambda_{it}\Theta_1p_{it}(\Theta_2),\\
\E{N_{it}^2\mid\Theta_1,\Theta_2}&=\lambda_{it}\Theta_1+\lambda_{it}^2\Theta_1^2,
&
\E{M_{it}^2\mid\Theta_1,\Theta_2}&=\lambda_{it}\Theta_1p_{it}(\Theta_2)
+\lambda_{it}^2\Theta_1^2p_{it}(\Theta_2)^2,
\end{aligned}
\]
and
\[
\begin{aligned}
\E{N_{it}M_{it}\mid\Theta_1,\Theta_2}
&=
p_{it}(\Theta_2)\E{N_{it}^2\mid\Theta_1,\Theta_2}
=
p_{it}(\Theta_2)\left(\lambda_{it}\Theta_1+\lambda_{it}^2\Theta_1^2\right),\\
\E{N_{is}M_{it}\mid\Theta_1,\Theta_2}
&=
\lambda_{is}\lambda_{it}\Theta_1^2p_{it}(\Theta_2),\\
\E{M_{is}M_{it}\mid\Theta_1,\Theta_2}
&=
\lambda_{is}\lambda_{it}\Theta_1^2p_{is}(\Theta_2)p_{it}(\Theta_2).
\end{aligned}
\]
Taking expectations, using the independence of \(\Theta_1\) and
\(\Theta_2\), and subtracting the products of the means proves the
result.
\end{proof}


%==========================================================================
% Section: Empirical analysis
%   Results produced by poisson_models (1).R (total claim counts) and
%   binomial_probit_models.R (extreme claim counts); see results_extreme/.
%==========================================================================
\section{Application to regional claim counts in Korea}
\label{sec.empirical}

In this section, we fit Model~\ref{model.thin} to a regional panel of
claim counts and evaluate its one-step-ahead predictions.

\subsection{Data}
\label{sec.emp.data}

The data set contains annual observations for \(I=17\) regions of
Korea over the years 2015--2024. For region \(i\) and year \(t\), the
total claim count \(N_{it}\) and the extreme claim count \(M_{it}\) are
recorded, and \(M_{it}\le N_{it}\) holds for all observations.
% TODO: state the definition of an extreme claim (e.g. the threshold) and the data source.
The covariate vector \(x_{it}\) consists of an intercept and the
logarithms of four regional characteristics: the recovery budget lagged
by one year (RecvBudget\_m1Y), the gross regional domestic product
(GRDP), the population, and the area.
% TODO: state the units of the covariates.
We use the years 2015--2023 as the training period and hold out 2024
for the evaluation of the predictions. The training period contains
\(153\) region-years, \(932\) claims, and \(200\) extreme claims. Since
\(8\) region-years have no claims, \(145\) region-years contribute to the
likelihood of the extreme claim part (Remark~\ref{rem.thin.estimation}).

Table~\ref{tab.emp.year} reports the annual totals. The total claim
count varies moderately over the years, whereas the share of extreme
claims varies considerably, from \(2.4\%\) in 2015 and \(4.1\%\) in 2021
to \(50.3\%\) in 2020. We return to this feature in
Section~\ref{sec.emp.discussion}.

\begin{table}[htbp]
\centering
\caption{Annual totals of the total and extreme claim counts over the
17 regions.}
\label{tab.emp.year}
\small
\begin{tabular}{lrrrrrrrrrr}
\toprule
Year & 2015 & 2016 & 2017 & 2018 & 2019 & 2020 & 2021 & 2022 & 2023 & 2024\\
\midrule
\(\sum_i N_{it}\) & 42 & 74 & 82 & 121 & 101 & 151 & 123 & 99 & 139 & 151\\
\(\sum_i M_{it}\) & 1 & 20 & 8 & 9 & 12 & 76 & 5 & 26 & 43 & 32\\
Share (\%) & 2.4 & 27.0 & 9.8 & 7.4 & 11.9 & 50.3 & 4.1 & 26.3 & 30.9 & 21.2\\
\bottomrule
\end{tabular}
\end{table}

\subsection{Estimation}
\label{sec.emp.estimation}

By Remark~\ref{rem.thin.estimation}, the maximum likelihood estimators
of \((\alpha,r)\) and \((\beta,\sigma)\) are obtained separately from
\(\ell^{(1)}\) and \(\ell^{(2)}\). The log-likelihood \(\ell^{(1)}\) is
available in closed form by \eqref{eq.thin.L1}. The log-likelihood
\(\ell^{(2)}\) is evaluated by the one-dimensional integral in
Remark~\ref{rem.thin.computation} using Gauss--Hermite quadrature with
\(100\) nodes; increasing the number of nodes to \(200\) changes
\(\ell^{(2)}\) by less than \(10^{-6}\), and adaptive numerical
integration agrees to within \(2\times10^{-4}\). Both log-likelihoods
are maximized by a quasi-Newton method with analytic gradients followed
by Newton iterations, and the standard errors are obtained from the
observed information matrix. The standard errors of the random effect
variances \(1/r\) and \(\sigma^2\) are obtained by the delta method.

As a benchmark, we fit the corresponding models without random effects,
namely the Poisson generalized linear model (GLM) for \(N_{it}\) and the
binomial probit GLM for \(M_{it}\) given \(N_{it}\). These correspond to
Model~\ref{model.thin} with \(\Theta_i^{(1)}\equiv1\) and
\(\Theta_i^{(2)}\equiv0\), so that the predictions are
\(\lambda_{i,t+1}\) and \(\Phi(\eta_{i,t+1})\), without experience
rating. We consider two specifications of \(x_{it}\): Specification A
uses all four covariates, and Specification B omits the recovery budget.
The random effect models are referred to as the generalized linear mixed
models (GLMMs).

Table~\ref{tab.emp.N} reports the estimates of the total claim count
part. The gamma random effect is supported by the likelihood ratio test,
whose \(p\)-value is computed from the boundary-corrected
\(\frac12\chi^2_1\) distribution, and by both information criteria. The
recovery budget is not significant in either model, and removing it
lowers both information criteria.

\begin{table}[htbp]
\centering
\caption{Estimates of the total claim count part \((\alpha,r)\), with
standard errors in parentheses. GLM: Poisson GLM; GLMM: Poisson-gamma
GLMM.}
\label{tab.emp.N}
\small
\begin{tabular}{lrrrr}
\toprule
& \multicolumn{2}{c}{Specification A} & \multicolumn{2}{c}{Specification B}\\
\cmidrule(lr){2-3}\cmidrule(lr){4-5}
& GLM & GLMM & GLM & GLMM\\
\midrule
Intercept & \(-3.192\) (1.024) & \(-0.096\) (2.492) & \(-2.559\) (0.737) & \(0.399\) (1.939)\\
log(RecvBudget\_m1Y) & \(0.060\) (0.067) & \(0.033\) (0.111) & -- & --\\
log(GRDP) & \(0.313\) (0.130) & \(1.120\) (0.432) & \(0.327\) (0.129) & \(1.179\) (0.390)\\
log(Population) & \(-0.117\) (0.134) & \(-0.910\) (0.415) & \(-0.115\) (0.133) & \(-0.958\) (0.389)\\
log(Area) & \(0.247\) (0.047) & \(0.230\) (0.088) & \(0.283\) (0.026) & \(0.247\) (0.068)\\
\(1/r\) & -- & \(0.091\) (0.066) & -- & \(0.100\) (0.064)\\
\midrule
Log-likelihood & \(-388.78\) & \(-382.69\) & \(-389.19\) & \(-382.73\)\\
AIC & \(787.55\) & \(777.38\) & \(786.38\) & \(775.47\)\\
BIC & \(802.71\) & \(795.56\) & \(798.50\) & \(790.62\)\\
LR test (\(p\)-value) & \multicolumn{2}{c}{\(12.18\) (\(0.0002\))} & \multicolumn{2}{c}{\(12.91\) (\(0.0002\))}\\
\bottomrule
\end{tabular}
\end{table}

Table~\ref{tab.emp.M} reports the estimates of the extreme claim part.
As for the total claim count, the Gaussian random effect is supported
by the likelihood ratio test in both specifications, with
\(\widehat\sigma^2=0.117\) in Specification A and
\(\widehat\sigma^2=0.058\) in Specification B, and the GLMM attains the
lower AIC and BIC. The recovery budget is again not significant, and
Specification B has the lower AIC and BIC among the GLMMs. In both
parts, the coefficients of log(GRDP) and log(Population) have opposite
signs and change considerably once the random effect is introduced.
These two covariates are strongly correlated across regions, and
population varies little within a region over time, so that only their
contrast, a GRDP per capita effect, is well identified. The coefficient
of log(Area) is the most stable across models; a larger area is
associated with a higher frequency and a higher share of extreme claims.
These associations are descriptive and should not be interpreted as
causal effects.

\begin{table}[htbp]
\centering
\caption{Estimates of the extreme claim part \((\beta,\sigma)\), with
standard errors in parentheses.}
\label{tab.emp.M}
\small
\begin{tabular}{lrrrr}
\toprule
& \multicolumn{2}{c}{Specification A} & \multicolumn{2}{c}{Specification B}\\
\cmidrule(lr){2-3}\cmidrule(lr){4-5}
& Probit GLM & Probit GLMM & Probit GLM & Probit GLMM\\
\midrule
Intercept & \(-0.220\) (1.735) & \(3.216\) (3.269) & \(0.246\) (1.105) & \(-0.098\) (1.808)\\
log(RecvBudget\_m1Y) & \(0.039\) (0.111) & \(-0.294\) (0.206) & -- & --\\
log(GRDP) & \(0.521\) (0.203) & \(1.033\) (0.448) & \(0.531\) (0.201) & \(0.733\) (0.305)\\
log(Population) & \(-0.650\) (0.209) & \(-0.968\) (0.420) & \(-0.652\) (0.208) & \(-0.779\) (0.316)\\
log(Area) & \(0.261\) (0.077) & \(0.444\) (0.148) & \(0.284\) (0.044) & \(0.269\) (0.063)\\
\(\sigma^2\) & -- & \(0.117\) (0.088) & -- & \(0.058\) (0.041)\\
\midrule
Log-likelihood & \(-310.05\) & \(-305.34\) & \(-310.12\) & \(-306.63\)\\
AIC & \(630.10\) & \(622.67\) & \(628.23\) & \(623.25\)\\
BIC & \(644.98\) & \(640.53\) & \(640.14\) & \(638.14\)\\
LR test (\(p\)-value) & \multicolumn{2}{c}{\(9.43\) (\(0.0011\))} & \multicolumn{2}{c}{\(6.98\) (\(0.0041\))}\\
Pearson \(\chi^2\)/df & \(4.32\) & -- & \(4.27\) & --\\
\bottomrule
\end{tabular}
\end{table}

\subsection{One-step-ahead prediction of 2024}
\label{sec.emp.prediction}

Using the training years as the history \(\mathcal H_{i,2023}\), we
compute the predictive probability \(\widehat p_{i,2024}\) of an extreme
claim from part iii of Theorem~\ref{thm.thin.pred}, and predict
\(M_{i,2024}\) in two ways:
\begin{itemize}
  \item[(a)] given the realized total claim count, by
  \(N_{i,2024}\,\widehat p_{i,2024}\) (part iv of
  Theorem~\ref{thm.thin.pred});
  \item[(b)] without it, by
  \(\E{N_{i,2024}\mid\mathcal H_{i,2023}}\,\widehat p_{i,2024}\)
  (part v of Theorem~\ref{thm.thin.pred}), where the first factor is
  given by part ii of Theorem~\ref{thm.thin.pred} with the estimates in
  Table~\ref{tab.emp.N}.
\end{itemize}
For the GLM benchmark, \(\widehat p_{i,2024}\) and
\(\E{N_{i,2024}\mid\mathcal H_{i,2023}}\) are replaced by
\(\Phi(\eta_{i,2024})\) and \(\lambda_{i,2024}\), respectively, with
the GLM estimates. The predictions are evaluated by the root mean
squared error (RMSE), the mean absolute error (MAE), and the deviance,
which is the binomial deviance in case (a) and the Poisson deviance in
case (b).

Table~\ref{tab.emp.crit} summarizes the results. In contrast to the
in-sample comparison, the GLMM does not improve the out-of-sample
predictions of the extreme claim counts: the GLM attains slightly lower
values of all three criteria in both cases and both specifications. In
case (b), the GLMM predicts a total of \(33.8\) and \(34.6\) extreme
claims, which is closer to the realized total of \(32\) than the GLM
totals of \(28.4\) and \(27.8\); the improvement in the total stems
from the total claim count part, for which the GLMM predicts the
realized total of \(151\) claims better than the GLM (\(133.4\) and
\(133.9\) against \(113.0\) and \(111.6\)).

\begin{table}[htbp]
\centering
\caption{Prediction criteria for the extreme claim counts in 2024
(\(17\) regions, realized total \(\sum_iM_{i,2024}=32\)).}
\label{tab.emp.crit}
\small
\begin{tabular}{llrrrr}
\toprule
& & \multicolumn{2}{c}{Specification A} & \multicolumn{2}{c}{Specification B}\\
\cmidrule(lr){3-4}\cmidrule(lr){5-6}
& Criterion & GLM & GLMM & GLM & GLMM\\
\midrule
(a) given \(N_{i,2024}\) & RMSE & \(2.341\) & \(2.524\) & \(2.323\) & \(2.569\)\\
& MAE & \(1.556\) & \(1.762\) & \(1.534\) & \(1.786\)\\
& Binomial deviance & \(37.10\) & \(43.30\) & \(36.66\) & \(43.83\)\\
& Predicted total & \(39.7\) & \(41.1\) & \(39.4\) & \(41.7\)\\
\midrule
(b) unconditional & RMSE & \(1.655\) & \(1.780\) & \(1.634\) & \(1.790\)\\
& MAE & \(1.170\) & \(1.304\) & \(1.157\) & \(1.326\)\\
& Poisson deviance & \(21.74\) & \(24.09\) & \(21.14\) & \(24.50\)\\
& Predicted total & \(28.4\) & \(33.8\) & \(27.8\) & \(34.6\)\\
\bottomrule
\end{tabular}
\end{table}

Table~\ref{tab.emp.region} reports the region-level results for
Specification B, ordered by the posterior mean of
\(\Theta_i^{(2)}\). The posterior means range from \(-0.30\) for
Chungnam to \(0.28\) for Gwangju, and the credibility adjustment moves
\(\widehat p_{i,2024}\) away from the a priori probability
\(\Phi\left(\eta_{i,2024}/\sqrt{1+\widehat\sigma^2}\right)\) by up to
\(11\) percentage points. The largest prediction errors of the GLMM
occur in regions whose share of extreme claims in 2024 departs from
their own history. For example, Gyeonggi and Chungnam had \(16\) of
\(119\) and \(14\) of \(70\) claims classified as extreme during the
training period, which lowers their predictive probabilities, but
\(7\) of \(13\) and \(7\) of \(12\) claims in 2024. Conversely, Jeonnam
and Gyeongnam had high shares during the training period but only
\(4\) of \(17\) and \(2\) of \(19\) in 2024.

\begin{table}[htbp]
\centering
\caption{Region-level results for Specification B: training totals,
posterior mean of \(\Theta_i^{(2)}\), a priori and predictive
probabilities of an extreme claim, and realized and predicted extreme
claim counts in 2024 given \(N_{i,2024}\) (case (a)).}
\label{tab.emp.region}
\footnotesize
\setlength{\tabcolsep}{4pt}
\begin{tabular}{lrrrrrrrrr}
\toprule
& \multicolumn{2}{c}{2015--2023} & & \multicolumn{3}{c}{Probability of extreme claim} & \multicolumn{3}{c}{2024}\\
\cmidrule(lr){2-3}\cmidrule(lr){5-7}\cmidrule(lr){8-10}
Region & \(\sum_tN_{it}\) & \(\sum_tM_{it}\) & \(\E{\Theta_i^{(2)}\mid\mathcal H}\) & A priori & GLM & GLMM & \(N_{i,2024}\) & \(M_{i,2024}\) & GLMM pred.\\
\midrule
Gwangju & 28 & 5 & \(0.278\) & 0.079 & 0.076 & 0.125 & 5 & 0 & 0.62\\
Jeonnam & 58 & 25 & \(0.214\) & 0.383 & 0.367 & 0.463 & 17 & 4 & 7.87\\
Busan & 46 & 5 & \(0.165\) & 0.085 & 0.075 & 0.110 & 5 & 0 & 0.55\\
Jeonbuk & 68 & 18 & \(0.162\) & 0.241 & 0.249 & 0.289 & 8 & 4 & 2.31\\
Seoul & 35 & 6 & \(0.145\) & 0.132 & 0.086 & 0.161 & 2 & 0 & 0.32\\
Gyeongnam & 64 & 19 & \(0.129\) & 0.301 & 0.284 & 0.344 & 19 & 2 & 6.53\\
Gyeongbuk & 94 & 38 & \(0.129\) & 0.393 & 0.376 & 0.440 & 9 & 3 & 3.96\\
Chungbuk & 70 & 22 & \(0.067\) & 0.333 & 0.317 & 0.354 & 10 & 3 & 3.54\\
Daejeon & 22 & 1 & \(-0.030\) & 0.082 & 0.079 & 0.077 & 6 & 1 & 0.46\\
Gangwon & 82 & 21 & \(-0.035\) & 0.333 & 0.345 & 0.317 & 22 & 1 & 6.97\\
Daegu & 22 & 1 & \(-0.036\) & 0.097 & 0.097 & 0.090 & 6 & 0 & 0.54\\
Sejong & 20 & 1 & \(-0.067\) & 0.100 & 0.110 & 0.088 & 2 & 0 & 0.18\\
Ulsan & 32 & 4 & \(-0.144\) & 0.260 & 0.220 & 0.214 & 6 & 0 & 1.28\\
Jeju & 45 & 3 & \(-0.146\) & 0.155 & 0.170 & 0.121 & 0 & 0 & 0.00\\
Incheon & 57 & 1 & \(-0.252\) & 0.115 & 0.101 & 0.072 & 9 & 0 & 0.65\\
Gyeonggi & 119 & 16 & \(-0.282\) & 0.280 & 0.227 & 0.191 & 13 & 7 & 2.48\\
Chungnam & 70 & 14 & \(-0.298\) & 0.392 & 0.354 & 0.283 & 12 & 7 & 3.39\\
\bottomrule
\end{tabular}
\end{table}

\subsection{Discussion}
\label{sec.emp.discussion}

The results indicate that the static random effects of
Model~\ref{model.thin} capture persistent differences between regions,
as shown by the likelihood ratio tests, but that the share of extreme
claims is dominated by shocks that are common within a year. Two
diagnostics support this interpretation. First, the Pearson statistic of
the probit GLM in Table~\ref{tab.emp.M} is about \(4.3\) times its
degrees of freedom, so the extreme claim counts are strongly
overdispersed relative to the binomial distribution. Second, adding
year fixed effects to the probit GLM of Specification B reduces the
deviance by \(156.9\) on \(8\) degrees of freedom, whereas adding region
fixed effects on top of them reduces it by a further \(41.0\) on \(16\)
degrees of freedom. Even with both sets of fixed effects, the Pearson
statistic remains about \(3.9\) times its degrees of freedom. This
suggests that the claims within a region-year are not conditionally
independent, as would be the case when several claims arise from the
same event. Under such shocks, the experience of a region over the
training period is only weakly informative about its share of extreme
claims in the next year, which explains why the credibility adjustment
of \(\widehat p_{i,2024}\) does not improve the predictions.

These findings are subject to two caveats. The evaluation is based on a
single hold-out year with \(17\) regions and \(32\) extreme claims, so
the differences in Table~\ref{tab.emp.crit} are small relative to their
sampling variability; a rolling-origin evaluation over several years
would give a more reliable comparison. Moreover, the conclusions concern
the static specification. Natural extensions of
Model~\ref{model.thin} are a common year effect, a random effect at the
region-year level to absorb the dependence among claims from the same
event, and dynamic random effects that discount older experience.



%\bibliographystyle{apalike}
%\bibliography{references}

\end{document}
